\documentclass[lstlisting]{hebdomon}
\addbibresource{refs.bib}

\title{Evaluating Cache-Replacement Policies Under a Mixed Workload}
\StudentName{P. Novak}
\date{\today}

\begin{document}

\maketitle
\tableofcontents

\Chapter{Objective}

This report evaluates three cache-replacement policies --- LRU
(Least Recently Used), LFU (Least Frequently Used) and CLOCK --- on a
simulated single-level cache under a workload that mixes a scanning
access pattern with a small set of hot keys. The goal is to identify
which policy best resists \emph{cache pollution} from the scan.

\Section{Background}

A pure scan (each key touched exactly once, in sequence) is the
textbook failure case for LRU: every scanned key evicts a
still-useful hot key, even though the scanned keys themselves are
never reused \parencite{belady1966}. LFU and CLOCK are both designed
to be more resistant to this, at the cost of extra per-entry
bookkeeping \parencite{orourke1990}.

\Chapter{Method}

\Section{Workload}

Each trial issues 200{,}000 accesses to a keyspace of 10{,}000 keys
split into two groups:

\begin{itemize}
    \item A \textbf{hot set} of 200 keys, together receiving 80\% of
    all accesses (Zipf-distributed within the set).
    \item A \textbf{scan set} of the remaining 9{,}800 keys, each
    touched at most twice, interleaved with the hot-set accesses at a
    fixed ratio of one scan access per four hot accesses.
\end{itemize}

The cache is sized at 500 entries (5\% of the keyspace, 2.5$\times$
the hot set), so a policy immune to scan pollution should converge to
a near-100\% hit rate on the hot set alone.

\Section{Implementation}

\begin{lstlisting}[language=Python,caption={CLOCK eviction core}]
class ClockCache:
    def __init__(self, capacity):
        self.capacity = capacity
        self.ref_bit = {}
        self.order = []
        self.hand = 0

    def access(self, key):
        if key in self.ref_bit:
            self.ref_bit[key] = 1
            return True  # hit
        self._evict_if_full()
        self.order.append(key)
        self.ref_bit[key] = 1
        return False  # miss
\end{lstlisting}

LRU was implemented with an ordered dict (O(1) move-to-front on
access); LFU with a min-heap keyed on access count, ties broken by
insertion order; CLOCK as shown above, with a circular scan and a
single reference bit per entry.

\Chapter{Results}

\begin{table}[h]
\centering
\caption{Overall and hot-set hit rate by policy}
\begin{tabular}{lrr}
\toprule
Policy & Overall hit rate & Hot-set-only hit rate \\
\midrule
LRU    & 61.2\% & 76.5\% \\
LFU    & 78.9\% & 98.1\% \\
CLOCK  & 74.4\% & 93.7\% \\
\bottomrule
\end{tabular}
\end{table}

LRU's hot-set hit rate of 76.5\% falls well short of the near-100\%
a scan-immune policy should achieve at this cache size, confirming
the expected pollution effect: each scan access evicts one hot key
whether or not it will ever be touched again. LFU comes closest to
immune, since a key touched only once or twice can never outrank a
hot key's accumulated count. CLOCK sits between the two --- its
single reference bit protects a hot key for one extra sweep past its
last access, but a long enough run of scan accesses between hot
touches still evicts it.

\Chapter{Discussion}

The results reproduce the standard LRU-scan-pollution story
quantitatively rather than just qualitatively: at this hot/scan
ratio and cache size, LRU loses roughly 1 in 4 hot-set requests to
scan-induced eviction that LFU recovers entirely. CLOCK's
middle-ground result matches its design intent as a cheap
approximation to LFU-like resistance without a full frequency
counter \parencite{orourke1990}, and the gap between CLOCK and LFU
here (4.4 percentage points on the hot set) is the price of that
approximation under this specific workload.

A caveat: LFU's advantage assumes the hot set is genuinely stable
over the measurement window. A workload whose hot set shifts over
time would erode LFU's advantage, since old high-count keys can
persist in cache long after they stop being useful --- a failure
mode this benchmark does not exercise.

\Chapter{Conclusion}

For a workload with a stable hot set contaminated by a one-shot
scan, LFU gives the best cache utilisation of the three policies
tested, with CLOCK a reasonable lower-overhead compromise and plain
LRU the clear loser.

\printbibliography

\end{document}
